\chapter{Randomness, paired controls and reproducibility}\label{ch:randomness}
``Use the same seed'' sounds like a complete reproducibility instruction.
It is not. Two implementations can start from the same seed and consume
draws in different orders, causing every later disturbance or choice to
diverge. Scientific pairing requires a more explicit account of randomness.

\section{Separate the sources of randomness}
The programme distinguishes nature or forcing draws from actor-selection
draws. If proposals are random, their stream also needs an explicit role.
A counter-addressed design can associate each draw with coordinates such
as stream, epoch, event and attempt, so an extra actor draw does not shift
future nature draws. This is a candidate implementation technique, not
a substitute for registering the actual stochastic law.

The generator, distribution transform, coordinate schema and retry behavior
all belong to the reproducibility contract. Independent seeds alone do not
prove independence of every downstream mechanism or prevent a helper from
reading a future draw.

\section{The physical-random comparison}
The proposed comparison uses the same declared physical domain and menu
construction but omits the EBU-capacity affordability filter. It is not an
EBU maximizer, a service optimizer or the historical full multi-edge P1C
arm. Those alternatives answer other questions.

Control parity means equal rules where equality was promised, not equal
realized actions. Once states differ, physical feasibility can make their
menus differ. A comparison must report that difference rather than force
one arm to perform an invalid action to keep a superficial pairing.

\section{Common raw draws versus common applied forcing}
Imagine a raw forcing draw proposes moving three units out of a cell.
One arm currently holds four there and another holds two. The same raw
proposal is feasible for one and infeasible for the other. Rejection,
clipping or resampling can yield different applied disturbances.

Thus shared raw randomness does not guarantee shared physical exposure.
A future protocol must declare the handling rule, record raw and applied
objects separately and determine how exposure differences affect the
comparison. Silently promising both forms of equality would be false.

This is not a reason to abandon paired designs. It is a reason to specify
what is paired. The scientific question may concern response under common
raw environmental opportunities, or response to identical admissible
physical inputs. They are different estimands when feasibility depends on
the evolving state.

\section{Restart without a new experiment}
A reproducible checkpoint must restore stock, balances, audit, parameter
version, pending event identity and random-stream coordinates together.
Restoring only the stock can leave an apparently identical physical world
with different future affordability or random choices.

Whether a restart is permitted after partial execution belongs to the
study's operational contract. A book cannot grant an extra scientific
attempt by calling it recovery. Historical Gate 1D-C illustrates why
execution attempts and completion sentinels are separately recorded.

\section*{Practice question}
Two arms have equal initial state and equal nature seed but use one shared
random generator for nature and actors. The EBU arm rejects more candidates
and consumes extra draws. Are future nature inputs paired? Not necessarily.
The rejection path has changed draw consumption. An independently addressed
nature stream would remove this particular coupling, while state-dependent
forcing feasibility would remain a separate issue.

\section*{Chapter recap}
Reproducibility concerns the exact stochastic process, not a seed label.
Fair comparison requires explicit stream ownership, exposure accounting
and restart semantics.
